\newcommand{\DoNotLoadEpstopdf}{}
\documentclass[11pt,reqno]{amsart}
\usepackage[T1]{fontenc}
\usepackage{lmodern,amsmath,amssymb,amsthm,mathtools,mathrsfs,microtype}
\usepackage[margin=1.1in]{geometry}
\usepackage{xr-hyper}
\usepackage[hidelinks]{hyperref}
\externaldocument[supp-]{build/supplement}[supplement.pdf]
\hypersetup{pdftitle={Collision records, arithmetic, and raw local inversion in a triangular Lorentz gas},pdfauthor={Qian Qi}}
\setlength{\emergencystretch}{2em}
\allowdisplaybreaks[2]
\numberwithin{equation}{section}
\newtheorem{theorem}{Theorem}[section]
\newtheorem{lemma}[theorem]{Lemma}
\newtheorem{proposition}[theorem]{Proposition}
\newtheorem{corollary}[theorem]{Corollary}
\theoremstyle{definition}
\newtheorem{definition}[theorem]{Definition}
\newtheorem{inputstatement}[theorem]{Input statement}
\theoremstyle{remark}\newtheorem{remark}[theorem]{Remark}
\newcommand{\R}{\mathbb R}\newcommand{\Z}{\mathbb Z}
\newcommand{\one}{\mathbf1}\newcommand{\dd}{\,\mathrm d}
\newcommand{\esssup}{\operatorname*{ess\,sup}}
\newcommand{\diver}{\operatorname{div}}
\newcommand{\TV}{\operatorname{TV}}
\title[Collision records and raw local inversion]{Collision records, arithmetic, and raw local inversion in a triangular Lorentz gas}
\author{Qian Qi}
\date{October 12, 2026. A2-DYN, revision 88}
\subjclass[2020]{37D50, 37A50, 60F05, 28A75}
\keywords{Sinai billiards, exact collision records, local limits, higher optical moments, weighted roof fibers, microscopic transport}
\begin{document}
\raggedbottom
\begin{abstract}
We study the exact-label pointwise local limit problem for actual
returns in a triangular Lorentz gas. A parameter-dependent refinement
of the physical homogeneous-strip argument gives capped incidence
localization with every exponent below $3/2$. Positive layer cake then
provides actual optical moments of every order $1<p<3/2$, without a
strong-space multiplier assertion. A complete segment trace estimate
converts these moments into a finite-count optical bound of order
$h^{1/10}$ on a family whose initial gradient and total residence
thresholds are $h^{1/8}$ and $h^{1/4}$. On its common survivors the
original weighted roof fibers satisfy a total-variation modulus of
order $h^{1/10}+\varepsilon^{-1}h^{1/4}$, uniformly at finite collision
count and at all incidence scales. True exits and the complementary
starting source remain explicitly charged. The original pointwise
objective is unchanged; the vanishing of its critical and zero-residence
source charge remains to be established.
\end{abstract}
\maketitle

\section{The unchanged physical problem}
Put $R_-=9/20$ and $R_+=47/100$, and let $R\in I=[R_-,R_+]$. Circular scatterers have centers in
$\Lambda=\Z(1,0)+\Z(1/2,\sqrt3/2)$. Write $T_R$ for the regular
collision map and $(\alpha,p)\in(\R/2\pi\Z)\times(-1,1)$ for its
coordinates; $p$ is the outgoing tangential velocity and
$\dd\nu=\dd\alpha\,\dd p/(4\pi)$. Physical flights satisfy
$\tau_-\le\tau_R\le\tau_+$, where $\tau_-=3/50$ and $\tau_+$
is the fixed finite-horizon bound.

The section is specified here to fix every event in the argument.
Put $h_0=1/200$. Let $Y_R$ be the union of the four open coordinate
squares of half-side $h_0$ centered at
\[
 (0,0),\quad(\pi/6,-1/2),\quad
 (-\theta_R,\sin\theta_R),\quad
 (\pi/3-\theta_R,\sin\theta_R).
\]
Here $\theta_R=2\arctan q_R$, and $q_R$ is the unique root in
$(49/100,53/100)$ of
\[
 8Rq(1+q^2)-4q(1-q^2)+\sqrt3(1-6q^2+q^4)=0.
\]
Set
\[
 U=\{|\alpha-\pi/6|<3/50,\ |p|<3/20\},\qquad
 Y_R^*=Y_R\cup U,\qquad c_*=\frac{91}{10000\pi}.
\]
The rectangles are disjoint and $\nu(Y_R^*)=c_*$. We always use
$\nu_R^*=c_*^{-1}\nu|_{Y_R^*}$, not a renormalization of a smaller
source.

At $m$ flights, let $L=L_{m,R}$ be their total length,
$K^{\rm c}_{m,R}$ their lattice displacement, and
$A_{m,R}=\sum_{j=0}^{m-1}\one_{Y_R^*}\circ T_R^j$. The event is
\begin{equation}\label{eq:v76-event}
 \Pi_{n,k,m,R}=Y_R^*\cap\{K^{\rm c}_{m,R}=k,\ A_{m,R}=n\}
                         \cap T_R^{-m}Y_R^*.
\end{equation}
The initial visit is counted; the terminal visit is tested separately.
On this event the actual $n$th return has collision count $m$ and
roof $L$. Denote its density by $p_{n,R}(k,m,t)$.

Write $c_j=\sqrt{1-p_j^2}$ and let $\Delta_j$ be the least clearance
of flight $j$ from any third disk, excluding its two incident disks.
Keep the original smooth function $0\le\Theta\le1$, zero on
$(-\infty,1]$ and one on $[2,\infty)$. With $q_0=47/53$,
$d_j=\min(j,m-j)$ and $d'_j=\min(j,m-1-j)$, set
\begin{align}
 H^\varepsilon&=
 \prod_{j=0}^{m}\Theta\!\left(\frac{c_j}{\varepsilon q_0^{d_j/4}}\right)
 \prod_{j=0}^{m-1}\Theta\!\left(\frac{\Delta_j}{\varepsilon q_0^{d'_j/4}}\right),
 \qquad W^\varepsilon=1-H^\varepsilon,\label{eq:v76-guard}\\
 b^\varepsilon(t)\dd t
   &=L_\#(W^\varepsilon\one_{\Pi_{n,k,m,R}}\nu_R^*).
 \label{eq:v76-source}
\end{align}
In particular first incidence and first clearance outside any selected
critical neighborhoods are still present. No section-distance factor
has been inserted into this guard.

\subsection{The arithmetic reference and the target}
For $y=(k_1,k_2,t,n)$ put
\[
 \mu_R=(0,0,\bar\tau_R,c_*),\quad
 Z=(y-m\mu_R)/\sqrt m,\quad
 \mathcal C_{m,R,n,k}(M)=\{t:|Z|\le M\}.
\]
We retain the finite spectral-peak reference, including peaks close to
but not on the unit circle. To specify its meaning, the collision
record observable is $(\kappa_1,\kappa_2,\tau,\one_{Y_R^*})$.
Let $\lambda_j(R,\vartheta)$ and $\mathsf P_j(R,\vartheta)$ be its
continued simple twisted eigenvalue and projection in each of the
fixed peak charts, and let $\chi_j$ be their frequency partition.
At the real maximum $a_j$ of $\log|\lambda_j|$ put
\[
 z_j=\lambda_j(R,a_j),\quad
 \mu_j=i^{-1}\nabla\log\lambda_j(R,a_j),\quad
 H_j=-\nabla^2\log\lambda_j(R,a_j),
\]
with $\operatorname{Re}H_j\ge aI_4$ for a common $a>0$. If
$\eta_R=\one_{Y_R^*}$ and $\ell$ denotes mass evaluation, put
$B_j=\chi_j(R,a_j)\ell(M_{\eta_R}\mathsf P_j(R,a_j)(\eta_R\nu))$.
Define
\begin{align}
 G_H(z)&=(2\pi)^{-4}\int_{\R^4}e^{-iv\cdot z-v^{\mathsf T}Hv/2}\dd v,
 \notag\\
 \mathcal L_{m,R}(y)&=\operatorname{Re}\left[
 c_*^{-1}\sum_j e^{-ia_j\cdot y}z_j^m B_j
      G_{H_j}\left((y-m\mu_j)/\sqrt m\right)\right].
 \label{eq:v76-reference}
\end{align}
Only the charts active at $R$ are summed. This is the same transition
reference as in the retained spectral theorem. The existence and
uniform estimates for these scalar data belong to that analytic
input, not to the transport argument below. No frequency branch or
arithmetic zero class is deleted.

The governing assertion remains
\begin{equation}\label{eq:v76-target}
 \sup_{R,\,1\le n\le m,\,k\in\Z^2}
 \esssup_{t\in\mathcal C_{m,R,n,k}(M)}
 |m^2p_{n,R}(k,m,t)-\mathcal L_{m,R}(k_1,k_2,t,n)|\longrightarrow0.
\end{equation}
Let $\mathcal E_M$ denote this error's collision limsup. An empty
central target contributes zero. The order of the uniformity is
$\sup_R\esssup_t$, so exceptional density sets may depend on $R$.

\subsection{Physical traces and a uniform microscopic envelope}
The sensitivity is the original additive function
$I_m=1+\sum_{i=1}^{m-1}q_0^{d_i}/c_i$. Write $z=|g|$ for the
endpoint-gradient norm and $\ell$ for the total residence of its
maximal regular exact-event segment. We use the common Borel coarea
kernels $\mu_t$. At fixed $m,R,n,k$ the identities hold outside a
null set in $t$, or a joint null set for product kernels; the uniform
norm is $\sup_R\esssup_t$, not a claim to one null set for all radii.

\begin{theorem}[Microscopic envelope on a shrinking initial family]
\label{thm:v88-main}
Set $a(h)=h^{1/8}$ and $\delta(h)=h^{1/4}$. On the original
reachable target fibers define
\[
 V_h^{\rm reg}(t)=\int_0^h\int
  \one_{\{\tau^-(y)>s,\ z(\Phi_{-s}y)\ge a(h),\ \ell(y)\ge\delta(h)\}}
                      \frac{I_m(y)}{z(y)}\,d\mu_{t+s}(y)\,ds.
\]
Under the explicitly retained inputs, there are fixed $K<\infty$
and $h_*>0$ such that
\begin{equation}\label{eq:v88-main-envelope}
 \sup_{m\ge2,R,n,k}m^2\|V_h^{\rm reg}\|_\infty
                              \le Kh^{1/10}\qquad(0<h<h_*).
\end{equation}
This includes all positive incidences and all values of $I_m$ on
the specified initial family. Both thresholds depend on the roof
width, not on the collision count, and the estimate holds uniformly
before any collision limit.
\end{theorem}
The proof is Lemma~\ref{lem:v88-trace} and
Theorem~\ref{thm:v88-expanded}. It uses the actual moment $I_m^{10/7}$
proved below without an inverse-gradient weight. The elementary trace inequality is proved in
full, including its integration over complete physical segments.

The complementary initial mass is not discarded. Define
\[
 R_h^\varepsilon(t)=\int W^\varepsilon
            \one_{\{z<a(h)\ \mathrm{or}\ \ell<\delta(h)\}}\,d\mu_t.
\]
Theorem~\ref{thm:v88-endpoint} identifies the original pointwise error
with the outer limit of the ordered microscopic height of this source.
The result does not assert that this last charge is zero, nor that
previous sufficient optical charges have vanished. The original
target, arithmetic reference and scalar-first downstream order remain
unchanged.

\section{Explicit analytic and geometric inputs}
We collect the indispensable inherited assertions in the exact form
used by the new proofs. Their source theorems and operator dependencies
are specified here and in Appendix~\ref{app:v83-kernels}. The new
moment and weighted-fiber estimates do not assume a surface spectral
insertion theorem or an unweighted pointwise local bound.

\begin{inputstatement}[Physical flow and source comparison]
\label{input:v83-geometric}
On the regular part $\Omega$ of each exact event, endpoint angles
$q=(\alpha_0,\alpha_m)$ give $F=L$, $g=\nabla_qF=R(-p_0,p_m)$,
$K=D_q^2F$ and $\vartheta=\arg g$. Put
\[
 \lambda_\circ=9\sqrt{19}/200,\qquad \Lambda_\circ=47/6,\qquad
 \omega=|\partial_{p_0}\alpha_m|^{-1},\qquad
 \chi=|\det D_q(F,\vartheta)|.
\]
The bounds $\lambda_\circ I\le K\le\Lambda_\circ I$ and
$\lambda_\circ\le\chi\le\Lambda_\circ$ hold uniformly.
The explicit field
\[
 X_q=\frac{K^{-1}g}{g^{\mathsf T}K^{-1}g}
\]
is stopped at true physical and section-decision boundaries, not at
chart seams. Its flow satisfies $XF=1$, $X\vartheta=0$ and
\begin{equation}\label{eq:v83-radial-input}
 2\lambda_\circ s\le |g(\Phi_sx)|^2-|g(x)|^2
                                      \le2\Lambda_\circ s.
\end{equation}
Its endpoint-curve length is at most $C\sqrt s$. For
$I_m=1+\sum_{i=1}^{m-1}q_0^{d_i}/c_i$ one has
$|X\omega|\le C\omega I_m/|g|$ on regular survivors. At a genuine
grazing exit some internal incidence tends to zero and $\omega$ tends
to zero. Forward interior critical termination cannot occur.

Every starting point with forward lifetime at most $r$ belongs to a
layer of one of the following forms, at an actual collision or flight:
\begin{equation}\label{eq:v83-layer-input}
 c_j\le C_Eq_0^{d_j}\sqrt r,\quad
 \Delta_j\le C_Eq_0^{d'_j}\sqrt r,\quad
 \operatorname{dist}(T_R^jx,\partial Y_R^*)\le C_Eq_0^{d_j}\sqrt r.
\end{equation}
The physical coarea covariance, proved in the appendix in the convention
used here, has Jacobian $J_h=(\omega(\Phi_hx)/\omega(x))\chi(x)/\chi(\Phi_hx)$.
Writing $E_h^\varepsilon$ for all weighted starting exit mass and
\[
 O_h^\varepsilon(t)=\int_{S_h(t)} W^\varepsilon(x)
            [1-\omega(\Phi_hx)/\omega(x)]_+\,d\mu_t(x),
\]
the complete source satisfies, for each fixed $r$ and almost every $t$,
\begin{equation}\label{eq:v83-shift-input}
 b^\varepsilon(t)\le2A_rb^\varepsilon(t)+\bar E_r^\varepsilon(t)
       +\bar O_r^\varepsilon(t)
       +C_\Theta(1+\varepsilon^{-1})\sqrt r\,A_rp(t),
 \qquad A_rf(t)=\frac1{2r}\int_{t-r}^{t+r}f(s)\,ds.
\end{equation}
\end{inputstatement}
These are the endpoint-field, all-contact motion, uniform cubic-action
and separated-volume results in retained modules 174, 177, 180--182.
The guard term in \eqref{eq:v83-shift-input} is the complete product
guard, not a protective subword source. The derivative and continuation
assertions are confined to regular physical trajectories.

\begin{inputstatement}[Complete local source estimates]
\label{input:v83-local}
For every fixed auxiliary band $B_*$, constants $C_{B_*}<\infty$ and
$\rho_{B_*}<1$ give the following bounds uniformly in radius, exact
labels, interval $J$ and mark $0\le j\le m$. Let $S_{R,s}$ be the
union of strips of width $Cs$ about the finitely many coordinate lines
supporting $\partial Y_R^*$. For $s$ small,
\begin{align}
 m^2\nu\{K_m^{\rm c}=k,A_m=n,L_m\in J,
                       T_R^jx\in S_{R,s}\}
 &\le Cs^{1/16}(|J|+B_*^{-1})
   +C_{B_*}(1+B_*|J|)m^2\rho_{B_*}^{m/2},\label{eq:v83-strip-input}\\
 m^2\int_Jb^\delta(t)\,dt
 &\le C\delta^{1/16}(|J|+B_*^{-1})
   +C_{B_*}(1+B_*|J|)m^3\rho_{B_*}^{m/2}.\label{eq:v83-guard-input}
\end{align}
The local extension to a fixed remaining compact range of $\delta$
is included. At auxiliary parameter one $W^1=1$, so $b^1=p$.

Put $z=|g|$, $\alpha=129/128$ and, for $0<\eta<1/2$, define the
actual nonnegative source measures
\begin{align}
 v^{\eta,b}(t)\,dt&=L_\#(I_m z^{-\eta}\one_{\{0<z<b\}}
                                       \one_\Pi\nu_R^*),\notag\\
 v^\eta(t)\,dt&=L_\#(I_m z^{-\eta}\one_\Pi\nu_R^*).
                                                   \label{eq:v83-fractional-input}
\end{align}
They are finite and absolutely continuous at every fixed count, and
for $0<b<b_1$,
\begin{equation}\label{eq:v83-fractional-finite-input}
 m^2\int_Jv^{\eta,b}(t)\,dt
 \le C_\eta b^{\alpha-\eta}(|J|+B_*^{-1})
  +C_{B_*,\eta}(1+B_*|J|)m^2b^{1/2-\eta}\rho_{B_*}^{m/4}.
\end{equation}
Adding the bounded-gradient part gives, for every fixed $L>0$,
\begin{equation}\label{eq:v83-fixed-width-input}
 \limsup_m\sup_{R,n,k}\sup_{|J|=L}m^2\int_Jv^\eta\le C_\eta L,
 \qquad
 \limsup_m\sup_{R,n,k}\sup_{|J|=L}m^2\int_Jv^{\eta,b}
                                      \le C_\eta b^{\alpha-\eta}L.
\end{equation}
Every auxiliary band is fixed before the collision limit and enlarged
only afterwards; leading constants are independent of that band and
of the previously fixed interval length.
\end{inputstatement}
Equation~\eqref{eq:v83-strip-input} is Supplement
Theorem~\ref{supp-thm:v48-marked-strip-local}. The physical source
bound \eqref{eq:v83-guard-input} is Supplement
Theorem~\ref{supp-thm:v49-physical-remainder-local}, in particular its
finite-window estimate. The fractional bounds are the positive
layer-cake conclusions of retained module 187, stated in full above.
The joint spectral and common-space assumptions used there are recorded
in the appendix. No estimate for an untruncated first inverse gradient
or for a singular roof convolution is part of this input.

\begin{inputstatement}[The scalar positive remainder]
\label{input:v83-scalar}
For large fixed $B$, with $\varepsilon(B)=A_0B^{-1/12}$,
\begin{equation}\label{eq:v76-remainder}
 m^2p-\mathcal L=m^2b^{\varepsilon(B)}+e_B,\qquad
 \limsup_m\sup_{R,n,k}\esssup_{\mathcal C(M)}|e_B|
                                      \le C_MB^{-1/192}.
\end{equation}
The source $b^\varepsilon$ is nonnegative and has a compactly supported
$BV$ density at fixed count. No collision-uniform variation bound is
assumed. Supplement Theorem~\ref{supp-thm:v51-positive-error} gives
this remainder with the same complete finite-peak reference.
\end{inputstatement}

\subsection{Maximal segments and the two existing charges}
The retained Borel parametrization, Supplement
Theorem~\ref{supp-v85archive:thm:v85-segment-map}, is a bijection
\[
 \Psi:\coprod_j\{(\theta,u):\theta\in E_j,\ a_j(\theta)<u<b_j(\theta)\}
                      \longrightarrow\Omega,
\]
with Borel inverse. Its exact density is
$A_j(u,\theta)=\omega(\Psi_j(\theta,u))/(4\pi c_*\chi(\Psi_j(\theta,u)))$.
Every segment remains in one regular exact event until a true boundary;
Borel chart seams are not exits. Nongrazing outgoing faces retain their
incident physical component and carry
$d\zeta_+=A_j(b_j-,\theta)\,d\theta$. The full angular variation is
$|\partial_u A_j|\le C A_j I_m/z$.
These are measure identities with the null-set convention already
specified, not bounds with an unknown number of words.

The true first-exit assignment is
\[
 \mathcal F_h(t)=\int_{\Gamma_{\rm ng}^+}
  \one_{\{0<b-t\le h,\ b-t<\tau^-(y)\}}\,d\zeta_+(y),\qquad
 \overline{\mathcal F}_r=r^{-1}\int_0^r\mathcal F_h\,dh.
\]
No lower gradient cutoff is imposed. Put $l=b-a$ and define
\[
 \mathcal Z_\delta(t)=\int_{\Gamma_{\rm ng}^+}
       \one_{\{0<l\le\delta\}}\one_{\{a<t<b\}}\,d\zeta_+,
 \qquad q_\delta(t)\,dt
  =F_\#(\one_{\Omega_+^{\rm ng}}\one_{\{l\le\delta\}}\bar\nu).
\]
The second object is the original collision volume restricted to the
same actual short segments, not a volume collar for a boundary.
For source families set
\[
 \mathcal H_M(f)=\limsup_m\sup_{R,n,k}m^2
            \esssup_{t\in\mathcal C_{m,R,n,k}(M)}f(t),
 \quad \mathfrak B_M=\lim_{\delta\downarrow0}\mathcal H_M(\mathcal Z_\delta),
 \quad \mathfrak Q_M=\lim_{\delta\downarrow0}\mathcal H_M(q_\delta).
\]
These limits are monotone extended nonnegative numbers.

For $0<\eta<1/2$, with $\kappa=(1-\eta)/2$, retain
\[
 \mathsf P_{\eta,h}v^\eta(t)=\int_0^h s^{-\kappa}v^\eta(t+s)\,ds,
 \qquad\mathfrak A_{\eta,M}
       =\lim_{h\downarrow0}\mathcal H_M(\mathsf P_{\eta,h}v^\eta).
\]
Supplement Theorems~\ref{supp-v85archive:thm:v85-cutoff-surface}
and \ref{supp-v85archive:thm:v85-boundary-charge} give
\begin{equation}\label{eq:v86-surface-input}
 \mathfrak B_M\le\mathcal H_M(\overline{\mathcal F}_r)
                         \le\mathfrak B_M+C_\eta r^{1/32}.
\end{equation}
Their band sum is finite before the collision limsup at each positive
residence cutoff. Its leading constant is cutoff independent; the
complete zero-residence tail is still inside $\mathcal H_M$ before
its cutoff tends to zero. The active offset correction to historical
module 189 remains Supplement equation
\eqref{supp-v83archive:eq:v83-dyadic-tail}; the historical inclusion is
explicitly marked superseded. None of these inputs asserts charge zero.

\subsection{The new finite-count theorems}
The broader moment below is unconditioned in the endpoint gradient;
it does not extend the inherited joint inverse-gradient moment range.
One finite roof-frequency band is fixed throughout its proof and its
microscopic applications. No claim of a band-independent leading constant
is needed for these finite-count envelopes.

For $a(h)=h^{1/8}$ and $\delta(h)=h^{1/4}$ the controlled family
includes every actual word with initial $z\ge a(h)$ and total residence
at least $\delta(h)$, at every positive internal incidence. Theorem
\ref{thm:v88-expanded} estimates its complete optical variation by
$Kh^{1/10}$. Theorem~\ref{thm:v88-fiber-stability} gives a modulus for
the actual weighted fibers, not only for a positive optical majorant.
Neither theorem declares true exits absent. The original exceptional
source is retained in Theorem~\ref{thm:v88-endpoint}.

\section{Optical moments below the homogeneous endpoint}
\label{sec:v88-moments}
The unconditioned incidence estimate permits a larger moment range than
its joint endpoint-localized counterpart. We keep these two estimates
separate. In particular, no improvement of the mixed inverse-gradient
moment is inferred from the theorem below.

Fix one admissible finite roof-frequency band $B_\dagger$. All constants
in this section may depend on this band and on the displayed exponents.
The band never depends on the collision count, an incidence threshold,
or a physical roof width.

\subsection{The retained spectral assumptions}
Write $Q=\mathscr T_{R,\xi}$, where $\xi$ is the real four-record
frequency, $M_s=M_{\one_{\{c<s\}}}$, and $\ell$ is mass evaluation.
We use the norm, multiplier, and physical-realization statements of
Supplement module 183, together with its fixed-band spectral input.
Specifically, the initial vector is $\one$, its physical iterates satisfy
$|Q^j\one|=1$, and
\begin{align}
 \|M_s f\|&\le C\|f\|,& |M_s f|_w&\le Cs^{1/8}\|f\|,
                                                    \label{eq:v88-multiplier}\\
 |M_sQ^j\one|_w&\le Cs^{3/2},&\|Q^j\|&\le C.
                                                    \label{eq:v88-real-weak}
\end{align}
The second estimate in \eqref{eq:v88-real-weak} is used on peak
charts. The first follows from homogeneous stable-curve length, in the
original invariant-density convention: incidence is $O(k^{-2})$ and
a homogeneous stable curve has length $O(k^{-3})$. No additional cosine
is inserted in the weak norm. Physical strong representatives and their
weak limits are the ones specified in the retained multiplier theorem.

On each peak chart the retained decomposition and inequalities are
\begin{align}
 Q^j&=\lambda^j\Pi+N^j,&\|N^j\|&\le C\rho^j,\notag\\
 \|Q^j f\|&\le C\sigma^j\|f\|+C(1+j)|f|_w,&
 |Q^j f|_w&\le C(1+j)|f|_w,                   \label{eq:v88-spectral}
\end{align}
with $\rho,\sigma<1$. The fixed-band complementary region has a
uniform spectral gap. The real peak powers have the four-dimensional
Gaussian integral bound $O(m^{-2})$. This paragraph states inherited
four-record inputs; it does not attribute them to a theorem for a
smaller record.

\begin{lemma}[Adjustable weak interpolation]\label{lem:v88-interpolation}
Let $0<\theta<1$ and $0<\epsilon_0<3/2$. A fixed refinement of the
peak charts gives
\begin{align}
 \|\Pi f\|+\sup_{a\ge0}|\ell(Q^a f)|
       &\le C\|f\|^{1-\theta}|f|_w^\theta,
                                        \label{eq:v88-interpolation}\\
 |M_s\Pi\one|_w&\le Cs^{3/2-\epsilon_0}.
                                        \label{eq:v88-projected}
\end{align}
All unit-modulus loci are retained. No normalization by $\Pi\one$ is
used, and the lemma remains valid when that vector vanishes.
\end{lemma}
\begin{proof}
Choose $\gamma>0$ with $1-\gamma>\theta$ and shrink the charts so that
$|\lambda|\ge r_*$ and
\[
 \frac{\log(1/r_*)}{\log(1/\sigma)}<\gamma,\qquad
 \frac{3\log(1/r_*)}{2\log(1/\rho)}<\epsilon_0.
\]
Uniform constants may first be chosen on larger fixed neighborhoods.
The remaining compact portions have a spectral gap and are still
estimated on the complement; they are not removed from the record or
from the final arithmetic reference.

Normalize $\|f\|=1$ and let $t=|f|_w<1/2$. Choose $L$ with
$\sigma^L\le t<\sigma^{L-1}$. Commutation and
\eqref{eq:v88-spectral} give
$\|\Pi f\|\le C(1+|\log t|)t^{1-\gamma}\le Ct^\theta$.
For $a\le L$ use the weak power inequality and continuity of mass;
for $a>L$ use the strong inequality at $L$ and bounded powers
thereafter. The cases $t=0$ and $t$ bounded below follow by passage to the limit
and boundedness, respectively. Scaling proves the first assertion.

For the second choose
$L=\lceil(3/2)\log(1/s)/\log(1/\rho)\rceil$ in the identity
\[
 M_s\Pi\one=\lambda^{-L}M_sQ^L\one-
                         \lambda^{-L}M_sN^L\one.
\]
Its weak norm is at most
$Cr_*^{-L}(s^{3/2}+s^{1/8}\rho^L)\le Cs^{3/2-\epsilon_0}$.
Its strong norm is bounded by \eqref{eq:v88-multiplier}. This uses
physical iterates, not an asserted modulus-one representative of a
continued eigenvector. The auxiliary integer $L$ is an interpolation
length, not a Fourier-band choice.
\end{proof}

\begin{theorem}[A capped marked estimate below $3/2$]
\label{thm:v88-incidence}
For every $0<\kappa<3/2$ and the fixed band $B_\dagger$, there
are $C_{\kappa,B_\dagger}<\infty$ and $\rho_\dagger<1$ such that
\begin{align}
 &m^2\nu\{K_m^{\rm c}=k,A_m=n,F\in J,c(T_R^j x)<s\}\notag\\
 &\quad\le C_{\kappa,B_\dagger}s^\kappa(|J|+B_\dagger^{-1})
  +C_{\kappa,B_\dagger}(1+B_\dagger|J|)m^2
                      \min\{\rho_\dagger^{m/2},s^2\}.
                                             \label{eq:v88-incidence}
\end{align}
The estimate holds for every finite $m\ge2$, every mark $0\le j\le m$,
interval $J$, exact coefficient and sufficiently small $s>0$.
No assertion is made at $\kappa=3/2$.
\end{theorem}
\begin{proof}
Choose $\theta,\epsilon_0$ so that
$(3/2-\epsilon_0)\theta>\kappa$. In the exact pairing
$\ell(Q^{m-j}M_sQ^j\one)$, at least one chronological block is
of length $m/2$. If the left block is long, expand it and apply
\eqref{eq:v88-interpolation} to $M_sQ^j\one$. The amplitude is
$O(s^{3\theta/2})$. If the right block is long, expand it and
apply \eqref{eq:v88-interpolation} to $M_s\Pi\one$, using
\eqref{eq:v88-projected}. The amplitude is
$O(s^{(3/2-\epsilon_0)\theta})$. Both are $O(s^\kappa)$.
The long-block remainders and the complementary frequency region are
$O(\rho_\dagger^{m/2})$, after changing fixed constants.

At every real frequency the same physical pairing also has absolute
value at most $\nu\{c<s\}=1-\sqrt{1-s^2}\le s^2$.
For nonnegative $A,D,V$, the two bounds $|z|\le A+D$ and
$|z|\le V$ imply $|z|\le A+\min(D,V)$. Thus the entire remainder
retains this physical cap. Apply the positive interval majorant with
roof-frequency support in the fixed band. Its Fourier supremum is
$O(|J|+B_\dagger^{-1})$. Gaussian integration of each long-block
peak term gives $m^{-2}$, whereas integration of the capped remainder
gives the second term of \eqref{eq:v88-incidence}. Exact torus
coefficient extraction retains $k,n$ separately. The finite refinement
has only changed the constants at the fixed band.
\end{proof}

\begin{theorem}[Actual complete optical moments below $3/2$]
\label{thm:v88-moment}
For every fixed $1<p<3/2$, the actual density
$u_p(t)\,dt=F_\#(I_m^p\one_\Pi\nu_R^*)$ is finite at every
collision count. On the one fixed band, for all intervals $J$,
\begin{align}
 m^2\int_Ju_p(t)\,dt\le C_{p,B_\dagger}(|J|+B_\dagger^{-1})
 +C_{p,B_\dagger}(1+B_\dagger|J|)m^2
                    \rho_\dagger^{m(2-p)/4}. \label{eq:v88-moment}
\end{align}
In particular, for $|J|\le1$,
\begin{equation}\label{eq:v88-uniform-energy}
       \sup_{m,R,n,k}m^2\int_Ju_p\le K_{p,B_\dagger}<\infty.
\end{equation}
No singular multiplier on the strong space is asserted. These are
unconditioned-gradient moments, not mixed $I_m^pz^{-\eta}$ bounds.
\end{theorem}
\begin{proof}
Choose $p<\kappa<3/2$. The positive incidence layer cake applied to
\eqref{eq:v88-incidence} has leading integral
$\int_0^{s_0}s^{\kappa-p-1}\,ds<\infty$. For the full remainder,
\[
 s_0^{-p}\min(D,s_0^2)+p\int_0^{s_0}s^{-p-1}\min(D,s^2)\,ds
      =\frac2{2-p}\min\{s_0^{2-p},D^{1-p/2}\}.
\]
This follows by splitting at $\sqrt D$ and holds for $p<2$.
Use $D=\rho_\dagger^{m/2}$ before removing any threshold. The
bounded-incidence part uses the retained unweighted fixed-band
concentration estimate.

Let $S=1+2q_0/(1-q_0)$. Weighted Jensen gives
$I_m^p\le S^{p-1}(1+\sum_iq_0^{d_i}c_i^{-p})$.
Sum the marked estimates with these geometric weights. Dropping the
endpoint tests on the positive upper comparison costs only $c_*^{-1}$
and leaves the exact coefficient unchanged; it does not redefine
$u_p$. This proves \eqref{eq:v88-moment}. The exponential remainder
has bounded $m^2$ supremum, proving \eqref{eq:v88-uniform-energy}.
Taking a containing roof interval proves finite mass at fixed count;
coarea on compact regular charts and positive exhaustion prove absolute
continuity in the common fiber version. Neither limit $p\uparrow3/2$
nor $B_\dagger\to\infty$ is used in the microscopic arguments below.
\end{proof}

\section{A larger shrinking family with a finite-count envelope}
\label{sec:v88-envelope}
The moment in Theorem~\ref{thm:v88-moment} improves the range of
initial geometries that can be controlled before the collision limit.
We write the trace argument for the full moment range, so its use does
not rely on substituting an exponent into the narrower historical
joint-moment theorem.

Fix $1<p<3/2$ and put $e_p=1-1/p$. At roof $t$ select every actual
maximal segment satisfying
\[
 a_j<t<b_j,\qquad z_j(t)\ge a,\qquad b_j-a_j\ge\delta.
\]
Denote this Borel family by $\mathcal E_t(a,\delta)$; $a$ is a
gradient threshold, not the endpoint roof $a_j$. Put
\[
 d=\min\{\delta/16,a^2/(32\Lambda_\circ),1/16\},\qquad
 J_t=(t-4d,t+4d),\qquad M_p(J)=\int_Ju_p.
\]
There is no restriction on internal incidences or on $I_m$.

\begin{lemma}[Complete traces for the enlarged moment range]
\label{lem:v88-trace}
For $0<a,\delta\le1$, in the common essential-roof versions,
\begin{equation}\label{eq:v88-trace}
 S_t:=\sum_j\int_{\mathcal E_t(a,\delta)}
    \sup_{t\le v<\min(t+d,b_j)}A_j(v,\theta)\,d\theta
                      \le C_p d^{-1}M_p(J_t).
\end{equation}
The displayed sum is over the disjoint segment parameter space.
\end{lemma}
\begin{proof}
One side of a selected segment has length at least $\delta/2\ge8d$.
Choose the anchor interval $(t-3d,t-2d)$ if the backward side is
long, and $(t+2d,t+3d)$ otherwise, with that Borel tie rule.
Every joining path from the anchor to a target stays in the same
physical segment and in $J_t$. Radial growth gives
$z(v)^2\ge a^2-8\Lambda_\circ d\ge3a^2/4$ on these paths.
The exact angular density therefore satisfies $|A'|\le(C/a)AI_m$.

For $f=A^{1/p}$ one has $|f'|\le(C/(pa))fI_m$ locally. Choose
an anchor point with $f\le(d^{-1}\int_{\rm anchor}A)^{1/p}$,
allowing an arbitrarily small error. The fundamental theorem of
calculus and H\"older on a joining interval of length at most $8d$
give
\[
 \sup f\le(d^{-1}\!\int_{J_t\cap(a_j,b_j)}A)^{1/p}
       +\frac C{pa}(8d)^{1-1/p}
                  (\int_{J_t\cap(a_j,b_j)}AI_m^p)^{1/p}.
\]
Since $d\le a$, raising to the power $p$ bounds the trace by
$C_pd^{-1}\int AI_m^p$. The Borel segment bijection identifies
the angular sum of this energy with a submeasure of $M_p(J_t)$,
with multiplicity one. Countable dense roof offsets give measurable
suprema; positive exhaustion gives the endpoint limits. This proves
\eqref{eq:v88-trace} for every $p>1$ for which the energy is finite.
\end{proof}

For $0<h\le d$ define the complete reachable variation source
\[
 \mathcal V_h^{a,\delta}(t)=\sum_j\int_{\mathcal E_t(a,\delta)}
    \int_0^{\min(h,b_j-t)}
       A_j(t+s,\theta)I_m(t+s,\theta)/z_j(t+s,\theta)
                                                     \,ds\,d\theta.
\]
\begin{theorem}[Trace-to-variation bound at every finite count]
\label{thm:v88-variation}
For every such $a,\delta,h$,
\begin{equation}\label{eq:v88-variation}
 \mathcal V_h^{a,\delta}(t)
 \le \frac{h^{e_p}}a S_t^{e_p}M_p((t,t+h))^{1/p}
 \le \frac{C_p}a(h/d)^{e_p}M_p(J_t).
\end{equation}
The complete sum includes all actual words of the fixed exact label.
\end{theorem}
\begin{proof}
On every selected forward segment $z(t+s)\ge a$. Apply H\"older
first in $s$ to $AI_m=A^{1-1/p}(AI_m^p)^{1/p}$ and then on the
disjoint angular parameter space. The summed first factor is exactly
$S_t$, not a maximum over words. This gives the first bound. Use
Lemma~\ref{lem:v88-trace} and $(t,t+h)\subset J_t$ for the second.
The integrals stop at the first true exit; no later re-entry is used.
\end{proof}

\begin{theorem}[An expanded finite-count microscopic family]
\label{thm:v88-expanded}
Set
\[
 p=10/7,\quad a(h)=h^{1/8},\quad\delta(h)=h^{1/4},\quad
 c_d=\min\{1/16,(32\Lambda_\circ)^{-1}\}.
\]
There is a fixed $h_*>0$ such that, for $0<h<h_*$,
$d=c_dh^{1/4}\ge h$, and
\begin{equation}\label{eq:v88-expanded}
 \sup_{m\ge2,R,n,k}m^2\|\mathcal V_h^{a(h),\delta(h)}\|_\infty
                              \le K_{B_\dagger}h^{1/10}.
\end{equation}
The same bound, up to a physical constant, pays surviving inverse-twist
loss and the original starting mass of genuine grazing first exits
on this family. The estimate is uniform at finite collision count.
\end{theorem}
\begin{proof}
Choose $h_*<\min\{1,c_d^{4/3}\}$. Insert
\eqref{eq:v88-uniform-energy} into \eqref{eq:v88-variation}.
The exponent is
\[
 -\frac18+\left(1-\frac7{10}\right)\left(1-\frac14\right)
                       =\frac1{10}.
\]
The constant is for the same single band at every $m$ and $h$.
On survivors negative optical variation pays the decrease of $\omega$;
at a genuine grazing exit it pays the full initial optical value
because the terminal value is zero. The angular normalization ratios
are uniformly bounded, so the same variation source bounds both
contributions on their disjoint initial sets. This does not assign
zero starting mass to a grazing exit.
\end{proof}

More generally, for $a=h^\xi$ and $\delta=h^\upsilon$, put
$b=\max\{2\xi,\upsilon\}$. The same theorem holds with exponent
$e_p(1-b)-\xi$ whenever $b<1$ and this exponent is positive.
For $\upsilon=2\xi$, every $\xi<1/5$ is admissible by a fixed
choice of $p<3/2$. This is the range delivered by this argument,
not a sharpness assertion for the billiard. The example above includes
states below both of the earlier thresholds $h^{1/128},h^{1/64}$.
It still leaves the complementary original starting source to be
estimated; cofinal cutoffs alone never prove its charge zero.

\section{Stability of the original weighted fibers}
\label{sec:v88-fibers}
The preceding estimate controls both signs of physical angular-density
variation. It therefore gives more than a bound on inverse-twist loss:
the original weighted source has a finite-count total-variation modulus
on its complete common-survivor family.

For $0<h\le d(a,\delta)$ let $\mathcal K_{t,h}$ be the members of
$\mathcal E_t(a,\delta)$ with $t+h<b_j$. On this same disjoint
parameter space define two measures
\[
 d\sigma_0=W^\varepsilon(\Psi_j(\theta,t))A_j(t,\theta)\,d\theta,
 \qquad
 d\sigma_h=W^\varepsilon(\Psi_j(\theta,t+h))A_j(t+h,\theta)\,d\theta.
\]
The second measure is the target fiber pulled back by the actual
segment correspondence. It is not the raw target restricted by new
labels, nor an independently renormalized source. We use the full
variation norm, so for these measures it is the integral of the
absolute difference of their densities, with no factor $1/2$.

\begin{theorem}[Original weighted roof-fiber stability]
\label{thm:v88-fiber-stability}
For almost every $(t,h)$ at finite count,
\begin{equation}\label{eq:v88-fiber-stability}
 \|\sigma_h-\sigma_0\|_{\rm var}
 \le C\mathcal V_h^{a,\delta}(t)
                 +K_\Theta\varepsilon^{-1}\sqrt h\,S_t(a,\delta).
\end{equation}
For the thresholds of Theorem~\ref{thm:v88-expanded},
\begin{equation}\label{eq:v88-fiber-finite}
 \sup_{m,R,n,k}m^2\esssup_t\|\sigma_h-\sigma_0\|_{\rm var}
 \le K_{B_\dagger,\Theta}
                   (h^{1/10}+\varepsilon^{-1}h^{1/4}).
\end{equation}
This holds for every fixed $\varepsilon>0$ and all sufficiently small
$h$ simultaneously. It controls every measurable test of absolute
value at most one under the same correspondence. Exited states are
not included in either measure and are charged separately below.
\end{theorem}
\begin{proof}
For each common survivor write $W_s=W^\varepsilon(\Psi(t+s))$.
Since $0\le W_s\le1$,
\[
 |W_hA_h-W_0A_0|\le|A_h-A_0|+A_h|W_h-W_0|.
\]
The full derivative estimate $|A'|\le CAI_m/z$ pays the first
term by $C\mathcal V_h^{a,\delta}$. It includes the change of
$\chi$ as well as $\omega$. The unchanged all-contact guard estimate
in Supplement Corollary~\ref{supp-v78archive:cor:v78-guard} gives
$|W_h-W_0|\le K_\Theta\varepsilon^{-1}\sqrt h$.
Summing $A_h$ over the common-survivor family costs at most $S_t$.
The identities use the common joint coarea kernels; the physical
Jacobian has been incorporated in $A_h$, not set equal to one.
This proves \eqref{eq:v88-fiber-stability}.

Lemma~\ref{lem:v88-trace} and \eqref{eq:v88-uniform-energy} give
$m^2S_t\le Kd^{-1}=Kc_d^{-1}h^{-1/4}$. Combine this with
\eqref{eq:v88-expanded}; multiplication by $\sqrt h$ leaves
$h^{1/4}$. Duality for signed measures proves the assertion for all
bounded measurable tests without selecting different null sets for
each test: equality of the two kernels is already an equality of
measures outside the common null set.
\end{proof}

\subsection{The complete original exceptional source}
Define, on the original fiber and with its original weight,
\begin{equation}\label{eq:v88-residual}
 R_h^{\varepsilon,*}(t)=\int W^\varepsilon
     \one_{\{z<h^{1/8}\ \mathrm{or}\ \ell<h^{1/4}\}}\,d\mu_t,
 \qquad \ell=\tau^-+\tau^+.
\end{equation}
Both alternatives include every endpoint type. The union is counted
once. Write $\mathfrak R_{\varepsilon,M}^*
=\lim_{h\downarrow0}\mathcal H_M(R_h^{\varepsilon,*})$.
Cofinality of the threshold pairs gives exactly
$\mathfrak R_{\varepsilon,M}^*=\mathfrak R_{\varepsilon,M}$
from Supplement Theorem~\ref{supp-v87archive:thm:v87-endpoint}.
Thus the larger controlled family does not by itself change or annul
the ultimate charge.

\begin{proposition}[A finite-count mass bound for the exceptional union]
\label{prop:v88-exceptional-mass}
On any fixed auxiliary band, with $\rho_*<1$ and fixed constants,
\begin{align}
 m^2\int_JR_h^{\varepsilon,*}(t)\,dt
 \le{}&C(h^{\alpha/8}+h^{1/128})(|J|+B_*^{-1})\notag\\
 &+C_{B_*}(1+B_*|J|)m^3\rho_*^{m/2},\qquad \alpha=129/128.
                                       \label{eq:v88-exceptional-mass}
\end{align}
Consequently its fixed-window collision-normalized mass is
$O(|J|h^{1/128})$, uniformly in $\varepsilon$. This assertion is a
mass bound, not a pointwise-height bound.
\end{proposition}
\begin{proof}
Drop $W^\varepsilon\le1$ only on the positive upper comparison.
The initial-gradient part uses the endpoint-only finite-count estimate
in Supplement equation~\eqref{supp-v86archive:eq:v86-endpoint-input}
at $\zeta=h^{1/8}$. If $\ell<h^{1/4}$, the forward lifetime is
less than $h^{1/4}$, regardless of the incoming or outgoing endpoint
type. The complete true-exit starting-mass bound of Supplement
Theorem~\ref{supp-v83archive:thm:v83-exit-local-mass} therefore applies
with width $h^{1/4}$ and contributes $h^{1/128}$. Add the two upper
bounds, which may overcount their intersection only on the upper
side. Enlarge $\rho_*<1$ to combine their remainders. The stated
fixed-window consequence takes the collision limit and then enlarges
the auxiliary band. No claim about the height of the union follows
from this operation.
\end{proof}

\begin{theorem}[The original pointwise budget with an expanded controlled family]
\label{thm:v88-endpoint}
For sufficiently small fixed $h>0$ and $\varepsilon>0$,
\begin{align}
 \mathcal H_M(b^\varepsilon)\le{}&
 \mathcal H_M(R_h^{\varepsilon,*})+C\varepsilon^{1/16}
       +Ch^{1/32}+Kh^{1/10}\notag\\
       &+C_\Theta(1+\varepsilon^{-1})\sqrt h.
                                              \label{eq:v88-endpoint}
\end{align}
With $\varepsilon(B)=A_0B^{-1/12}$, it follows that
\begin{equation}\label{eq:v88-squeeze}
 \mathfrak R_{\varepsilon(B),M}^*\le\mathcal E_M+C_MB^{-1/192},
 \qquad
 \mathcal E_M\le\mathfrak R_{\varepsilon(B),M}^*+C_MB^{-1/192}.
\end{equation}
The remaining assertion for the original scalar theorem is
$\lim_{B\to\infty}\mathfrak R_{\varepsilon(B),M}^*=0$ for every $M$.
It is not proved by the mass or common-survivor stability theorem.
\end{theorem}
\begin{proof}
Split the initial fiber before estimating any exit. On the exceptional
union, survivor and exit are disjoint; their positive defect is at most
the original weight, hence at most $R_h^{\varepsilon,*}$ in total.
On its complement, optical variation pays survivors and genuine grazing
exits by \eqref{eq:v88-expanded}. For a nongrazing first exit,
$\omega(t)\le\omega(b-)+\int_t^b[-\omega']_+$.
The endpoint term, with its physical angular normalization, is bounded
by a fixed multiple of the same outgoing first-exit assignment on
segments of total residence at least $\delta(h)=h^{1/4}$.
Its ordered height is at most $Ch^{1/32}$ by the retained positive
residence theorem, with this positive cutoff fixed before the collision
limit. The pullback Jacobian cost is paid by the variation term.

Insert these bounds into the complete source-shift inequality
\eqref{eq:v83-shift-input}. The original fixed-width bad source and
unweighted source estimates pay its average and guard terms. This proves
\eqref{eq:v88-endpoint}. Here the sharper ordered guard estimate is
used; \eqref{eq:v88-fiber-finite} is the stronger finite-count statement
for arbitrary measurable tests on the common-survivor family.

Now let $h$ decrease after $\mathcal H_M$, keeping $\varepsilon$
fixed. Since $R_h^{\varepsilon,*}\le b^\varepsilon$, the resulting
height differs from its residual charge by at most
$C\varepsilon^{1/16}$. The unchanged positive scalar remainder gives
\eqref{eq:v88-squeeze} by the triangle inequality. All arithmetic
classes remain unnormalized. The prior sufficient charge pairs remain
in the complete archive; this direct-source equivalence neither proves
those larger charges zero nor promotes downstream laws before scalar
closure.
\end{proof}


\section{The unchanged pointwise objective}
The new finite-count results quantify a larger controlled family and
the stability of its original weighted fibers. They are not a theorem
at zero initial gradient or zero total residence. The mass estimate
\eqref{eq:v88-exceptional-mass} likewise does not bound that source's
normalized essential height. The remaining original-source assertion
is the zero limit in Theorem~\ref{thm:v88-endpoint}.

The prior sufficient pairs of fractional or reachable optical charge
and zero-residence surface or volume charge remain valid in the complete
archive. No such larger charge is declared zero by a theorem on the
common-survivor family. Once the scalar pointwise endpoint is proved,
Supplement Theorem~\ref{supp-thm:v73-scalar-path-equality} gives the
same-roof path numerator. Reference normalization is restricted to
positive arithmetic classes; the clock and likelihood hypotheses remain
separate, and zero classes are never divided by.

\appendix
% Full inherited normalization and kernel proofs, with the new trace interface.
% Retained normalization proof, with three explicit references redirected to the archive.
\section{Physical kernels, first-exit traces and analytic dependencies}
\label{app:v83-kernels}
This appendix supplies the measure-theoretic interfaces used by the new
arguments in one convention. It does not deduce a surface estimate from a
volume multiplier theorem. The geometric and spectral estimates whose
proofs are retained in the technical supplement are listed separately.

\subsection{Normalization and common kernels}
On a regular word use endpoint angles $q=(\alpha_0,\alpha_m)$ and
oriented internal arclengths. Endpoint first variation gives
$F_{\alpha_0}=-Rp_0$ and $F_{\alpha_m}=Rp_m$. Consequently
\[
 \omega=|\partial_{\alpha_m}p_0|=|F_{\alpha_0\alpha_m}|/R,
 \qquad d\bar\nu=\frac{\omega}{4\pi c_*}\,dq.
\]
In the optical Schur complement the internal Hessian is $D_cAD_c$,
where $A$ has diagonal $t_{i-1}+t_i+2/(Rc_i)$ and off-diagonal
$-t_i$. The two angle-to-arclength factors are $R^2$; division by
$R$ in the first variation gives $\omega=Rc_0c_m/y_m$. For one
flight $y_1=\ell_0$, and for two flights
$y_2=\ell_0+\ell_1+2\ell_0\ell_1/(Rc_1)$. Changes of tangent
orientation do not change these absolute identities.

Away from $g=0$, the coordinates $(F,\vartheta)$ have Jacobian
$\chi=e^\perp{}^{\mathsf T}Ke^\perp$ in absolute value, where
$e=g/|g|$. Thus the angular expression for the original fiber is
\begin{equation}\label{eq:v83-appendix-fiber}
 d\mu_t=\frac{\omega(t,\vartheta)}{4\pi c_*\chi(t,\vartheta)}
          \,d\vartheta.
\end{equation}
The sign of an oriented endpoint chart is irrelevant; all source
Jacobians in this appendix are positive. The original source already
contains its initial and terminal tests and is normalized only by $c_*$. 

\begin{proposition}[Common graph covariance]
\label{prop:v83-common-kernel}
At fixed $m,R,n,k$, the regular exact event has jointly measurable
coarea kernels for which, for almost every $(t,h)$ and every nonnegative
Borel graph weight $f$,
\begin{equation}\label{eq:v83-common-covariance}
 \int_{S_h(t)}f(x,h,\Phi_hx)J_h(x)\,d\mu_t(x)
 =\int_{\Phi_hS_h(t)}f(\Phi_{-h}y,h,y)\,d\mu_{t+h}(y).
\end{equation}
The same versions can be used for countably many source restrictions,
positive truncations and their increasing limits. Here $S_h(t)$ is the
set of source points whose forward lifetime is greater than $h$.
\end{proposition}
\begin{proof}
At fixed count the regular physical words are disjoint open pieces
of the collision cylinder. Their true boundaries and the removed
critical points have zero source area. Take a countable cover by
precompact flow boxes with lifted angular coordinate. Assign each
point to its first containing box to obtain a Borel disjointization.
This assignment is used only for integration: its seams do not stop
the physical flow.

On the open survivor graph the map $(x,h)\mapsto(\Phi_hx,h)$ is
injective, smooth and has physical Jacobian $J_h$. Its inverse on
its image is the backward flow. Since $F(\Phi_hx)=F(x)+h$, ordinary
change of variables followed by disintegration over $(F,h)$ gives
\eqref{eq:v83-common-covariance} on a countable determining algebra
of Borel sets. Remove the union of the resulting null sets. Uniqueness
of the kernel measures and monotone approximation extend equality to
all nonnegative Borel weights outside this one joint null set.
Countable disjoint summation gives the complete physical formula.

In endpoint coordinates the field is translation in $F$ at fixed
$\vartheta$, so its endpoint Jacobian is $\chi(x)/\chi(\Phi_hx)$.
Multiplication by the physical density ratio gives
$J_h=(\omega(\Phi_hx)/\omega(x))\chi(x)/\chi(\Phi_hx)$, exactly
the normalization in \eqref{eq:v83-common-covariance}. Nonnegative
graph selectors accommodate first-exit restrictions depending on
$(x,h)$ without new incompatible versions. Monotone convergence at
fixed count permits the actual fractional insertions. No null set
independent of the uncountable radius parameter is required: the
ordered norm is $\sup_R\esssup_t$.
\end{proof}

\subsection{True outgoing faces and exhaustion}
On a nongrazing limiting word all incidence limits are positive, so
the stationary reflection equations extend smoothly to the boundary.
This continuation is not asserted to remain the physical first-hit
trajectory beyond that boundary. On a transverse outgoing arc write
$F=b(\vartheta)$. Contraction of the area form
$\omega(4\pi c_*\chi)^{-1}dF\,d\vartheta$ with $\partial_F$
gives the positive flux in \eqref{supp-v85archive:eq:v83-surface-flux}. A first-exit
point is always counted together with its incident one-sided physical
component. Lower-dimensional and tangent strata have zero transverse
flux. Section decisions, competing obstacles and simultaneous failures
are not suppressed by this description. An exit with a vanishing
internal incidence belongs to the grazing class even when another
condition fails there simultaneously; the remaining class is nongrazing.

\begin{proposition}[Initial mass and actual first-exit flux]
\label{prop:v83-first-exit-kernel}
For a nongrazing outgoing face, let $s=F(y)-t$ and $x=\Phi_{-s}y$.
On $0<s\le h$, $s<\tau^-(y)$, the contribution to initial weighted
exit mass is
\begin{equation}\label{eq:v83-inverse-jacobian-exit}
                 \frac{W^\varepsilon(x)}{J_s(x)}\,d\zeta(y).
\end{equation}
The boundary source \eqref{supp-v85archive:eq:v83-untruncated-flux} has this same
first-exit selection with weight one on the boundary, not with the
factor in \eqref{eq:v83-inverse-jacobian-exit}. Both are defined on
common almost-everywhere starting-roof versions.
\end{proposition}
\begin{proof}
In one angular flow box put
$a(u,\vartheta)=\omega(u,\vartheta)/(4\pi c_*\chi(u,\vartheta))$.
The initial fiber mass is $W(x)a(t,\vartheta)d\vartheta$, whereas
the outgoing flux is $a(b(\vartheta),\vartheta)d\vartheta$.
The ratio of the latter density to the former unweighted density is
$J_s(x)$. This proves \eqref{eq:v83-inverse-jacobian-exit}. The
backward-lifetime condition says precisely that the segment between
the starting roof and this boundary has remained in that one-sided
event; a later re-entry does not satisfy that condition.

Use precompact regular flow boxes, with rational rectangles in
$(F,\vartheta)$, and increasing finite unions whose interiors cover
the regular event. On overlaps within a physical component their
coordinates agree up to a lifted angular constant and can be refined
into finitely many pieces. Roof-constant faces are transverse;
angle-constant faces have zero flux. The outside source plus the
first-exit source for each such domain is the integral of the
indicator that the whole forward segment is not contained in the
domain. These indicators decrease to the indicator of a true exit.
Every compact surviving segment eventually lies in an exhaustion
domain. The usual endpoint convention at lifetime exactly $h$ does
not change the joint almost-everywhere formula.

For actual nongrazing faces, exhaust the regular one-sided arcs and
their backward segments by such compact pieces. Disjoint assignment
by the incident component and first-exit graph removes chart
multiplicity. Apply the local formula and nonnegative convergence
at fixed count. These operations also justify the variation payments
in Theorem~\ref{supp-v85archive:thm:v83-two-sided-flux}. No surface limit is commuted
with the collision-normalized supremum.
\end{proof}

\subsection{The indispensable analytic chain}
Input~\ref{input:v83-local} records the complete estimates used by this
article, rather than requiring a reconstruction from historical status
files. Its common-space operator dependencies are the collision norm
definitions (Supplement equations~\eqref{supp-eq:v35-weak-norm}--
\eqref{supp-eq:v35-unstable-norm}), the full-record fixed-band
Doeblin--Fortet and spectral estimates (Supplement
Theorem~\ref{supp-thm:v40-polynomial-LY}), and the finite moving-peak
local inversion. The proof uses bounded transverse multipliers, not a
bounded strong operator of multiplication by $c^{-1}$ or $|g|^{-1}$.

For specificity, the joint endpoint/incidence estimate has leading
factor $b^{129/128}s^{17/16}$ and finite-count cap
$\min\{b\rho^{m/2},s^2\}$. The endpoint test is the particular
vector $P_b\mathbf1$ of strong norm $O(b^{7/8})$; that estimate is
not a small operator norm on arbitrary vectors. Each continued
rank-one peak is normalized using a nonzero smooth test against a
physical peripheral vector, never by division by $\Pi\mathbf1$.
The homogeneous trace estimates pass through bounded finite Cesaro
approximants in the common strong space. These conventions and the
two chronological characteristic pairings are proved in the retained
module 186, not assumed for arbitrary $L^\infty$ representatives.

Positive incidence layer cake integrates the cap first and yields
a critical-volume remainder $b^{1/2}\rho^{m/4}$. A second layer cake
against $b^{-1-\eta}$ converges at finite count only for
$\eta<1/2$, giving the actual $v^{\eta,b}$ estimate stated in
Input~\ref{input:v83-local}. The leading exponent is $129/128-\eta$.
The first inverse-gradient remainder instead depends on its lower
cutoff and is not used to remove that cutoff here. For the section
layer only the all-depth marked-strip theorem is needed; for the
incidence and clearance layers the original physical-guard local
bound is used. All of these sources keep the forward occupation
coefficient. Neither signed interface cancellation nor deletion of
an arithmetic peak enters the new estimates.

Thus the new potential and surface theorems depend on explicitly
stated volume estimates, but the surface source is bounded only after
its actual first-exit graph and optical variation have been treated.
The retained spectral and billiard estimates remain available for
specialist verification. This appendix and native source compilation
are not independent human or formal proof certification.


\subsection{The new trace interface}
The new proof uses the exact segment density and the inequality
$|\partial_u A|\le CAI_m/z$, including both the optical and angular
factors. Applying the fundamental theorem of calculus to $A^{1/p}$
requires only local absolute continuity and the actual integral of
$AI_m^p$; it is not an anisotropic multiplier assertion. The reference
interval is chosen on the longer side of each actual segment, with a
Borel tie rule. Fubini then sums its energy with multiplicity one.
The resulting finite-count bound permits the two initial thresholds
to shrink with roof width. It does not extend through the unestimated
critical/zero-residence source, which is retained in the final charge.

\subsection{The enlarged moment and weighted-fiber interface}
The norms and physical multiplier realization used in
\eqref{eq:v88-multiplier}--\eqref{eq:v88-spectral} are precisely
Supplement equations~\eqref{supp-eq:v35-weak-norm}--\eqref{supp-eq:v35-unstable-norm},
Supplement Proposition~\ref{supp-prop:v49-physical-multiplier}, and
Supplement equations~\eqref{supp-v80archive:eq:v80-physical-modulus}
and \eqref{supp-v80archive:eq:v80-spectral-input}. The fixed-band
refinement used here is quantitative in the chosen moment exponent;
it uses the physical real-power length bound, not traces of an arbitrary
bounded measurable representative. Continued projections are handled by
finite physical powers and a remainder, even when $\Pi\one=0$.

The resulting $I_m^p$ moment has no inverse-gradient insertion. Its
constants may depend on the one fixed band. This is sufficient for the
finite-count supremum in \eqref{eq:v88-expanded}; an additional
band-limit claim would require additional uniformity and is not made.
The original guard modulus is Supplement
Corollary~\ref{supp-v78archive:cor:v78-guard}. The weighted-fiber
measures use one common segment correspondence and the full variation
norm. They do not compare separately selected labels or neglect exited
mass. The earlier trace interface remains in the retained source with
its original, narrower moment range.

\begingroup
\small
\enlargethispage{2\baselineskip}
\begin{thebibliography}{99}
\bibitem{AFP} L. Ambrosio, N. Fusco and D. Pallara,
\emph{Functions of Bounded Variation and Free Discontinuity Problems},
Oxford Mathematical Monographs, Clarendon Press, Oxford, 2000.
\bibitem{CMInt} R. Cluckers and D. J. Miller,
\emph{Stability under integration of sums of products of real globally
subanalytic functions and their logarithms}, Duke Math. J.
\textbf{156} (2011), no. 2, 311--348; arXiv:0911.4373v1,
Theorem 1.3.
\bibitem{Coste} M. Coste, \emph{An introduction to o-minimal geometry},
lecture notes, November 1999; Dottorato di Ricerca in Matematica,
Universit\`a di Pisa, Istituti Editoriali e Poligrafici Internazionali, 2000.
\end{thebibliography}

\endgroup
\end{document}
